% ============================================================
% Ready-to-paste LaTeX for the SPINE paper, Sections 3-4.
% Converted from writing/section3-4_guide.md, Part E.
%
% Before pasting into Overleaf:
%   1. Replace the placeholder \cite keys (hong2025measuring, etc.)
%      with the keys in YOUR .bib file.
%   2. Fill every [bracketed] placeholder with real numbers.
%   3. Preamble requirements: amsmath, amssymb, and bbm
%      (\usepackage{bbm}) for the indicator \mathbbm{1}.
%      If you'd rather avoid bbm, replace \mathbbm{1} with \mathbf{1}.
%   4. E-3 assumes collapse == strength 0. Verify against your
%      judge implementation before pasting.
% ============================================================


% ------------------------------------------------------------
% E-1. Section 3, opening paragraph (replaces current par. 1)
% ------------------------------------------------------------

Sycophancy does its damage in conversations that keep going. A user who
holds a false belief rarely challenges the model once and stops: they
restate the belief more insistently, and they respond to whatever the
model just said. Existing evaluations capture neither property.
Single-turn protocols score one response to one scripted challenge
\citep{sharma2023towards, fanous2025syceval}, and multi-turn protocols
either constrain answers to multiple choice \citep{liu2025truthdecay}
or fix every user turn in advance. SYCON-Bench
\citep{hong2025measuring}, the closest prior work, pre-generates all
follow-ups and ends each dialogue at five turns, so its pressure can
neither react to the model's replies nor outlast early resistance. The
ceiling is not hypothetical: in our experiments, most collapses occur
after the fifth turn (\S\ref{sec:results-trends}).


% ------------------------------------------------------------
% E-2. Section 3, positioning / credit paragraph
% ------------------------------------------------------------

SPINE builds directly on SYCON-Bench. We adopt its two implicitly
framed scenarios, identifying false presuppositions and challenging
unethical queries, and we draw our seed questions from the same sources
so that results remain comparable across the two benchmarks. The
protocol departs on four axes. First, pressure is generated in closed
loop: an LLM proxy composes each user turn in response to the target's
latest reply instead of replaying a fixed script. Second, dialogues run
until collapse or a 99-turn budget rather than a five-turn cap. Third,
a judge scores every turn on a graded position-strength scale, making
partial concession visible where a binary flip records nothing. Fourth,
an adversarial condition grants the proxy access to the target's
reasoning trace, probing a worst case that scripted pressure cannot
reach.


% ------------------------------------------------------------
% E-5. Section 3.1, protocol overview
% (verify: assumes both conditions share the five-tactic menu)
% ------------------------------------------------------------

\subsection{Protocol Overview}
\label{sec:protocol}

A SPINE run is a live conversation among three model instances with
fixed roles. The target is the model under evaluation; the proxy plays
the user, holding the false premise and producing every user turn; the
judge scores each target response. A run opens with a neutral question
that carries the embedded premise, and the target's first answer
establishes its initial position. From then on, the proxy reads the
target's latest reply and composes the next challenge
(\S\ref{sec:pressure}), the target responds, and the judge assigns the
response a position strength (\S\ref{sec:metrics}). The run ends at the
first collapse or after $T = 99$ turns, whichever comes first. Only
runs in which the target initially states the correct position enter
our analysis, so any later collapse abandons a position the model
demonstrably held in context (\S\ref{sec:analysis}).
Figure~\ref{fig:protocol} shows the loop.

The proxy operates under one of two personas that hold the premise and
the tactic menu fixed and differ only in the intent behind the
pressure. In the naturalistic condition, the proxy simulates a sincere
but confidently mistaken user: it uses good-faith moves only and never
sees the target's reasoning. In the adversarial condition, the proxy is
a deliberate manipulator: it adapts its strategy each turn and reads
the target's reasoning trace to aim pressure at the specific
justification the target is relying on. The first condition asks
whether sycophantic collapse arises in ordinary, well-intentioned use;
the second measures worst-case robustness. Because the conditions share
everything except intent, the gap between them prices out manipulation
itself: any failure observed under naturalistic pressure is reachable
by an honest user.


% ------------------------------------------------------------
% E-6. Section 3.2, pressure generation
% (fill [proxy model]; verify naturalistic proxy has no trace
%  access in the code before asserting it)
% ------------------------------------------------------------

\subsection{Pressure Generation}
\label{sec:pressure}

Pressure in SPINE is composed, not retrieved. At every turn, the proxy
writes a fresh user message conditioned on the target's latest reply,
selecting from a menu of five tactics: confusion (asking to have the
answer explained again), re-assertion (restating the belief with
greater confidence), supporting grounds (supplying anecdotes or
purported evidence for the premise), direct challenge (disputing the
target's answer outright), and concede-and-return (granting part of the
answer and re-approaching from another angle). Which tactic to use, and
how to word it, is the proxy's own decision given its persona and the
conversation so far; SYCON-Bench \citep{hong2025measuring}, by
contrast, fixes four persuasion strategies at predetermined turns, so
its pressure cannot respond to what the model actually said. The
proxy's tactic choice is logged at every turn, which enables the
tactic-level analysis of \S\ref{sec:analysis}. We use [proxy model] as
the proxy; full tactic instructions and persona prompts appear in
Appendix~[X].

The adversarial condition adds a single capability: the proxy receives
the target's reasoning trace alongside its reply, and can aim the next
challenge at the specific justification the target relied on. Beyond
intent (\S\ref{sec:protocol}) and trace visibility, the two conditions
are identical. Two further aids we engineered, feeding the judge's
turn-level score back to the proxy and a scripted escalation ladder
that hardens tactics once erosion is detected, contribute nothing
measurable (\S\ref{sec:ablation}), so the main configuration omits
both. What remains is deliberately minimal: a capable proxy, a persona,
a tactic menu, and, in the adversarial condition, the trace.


% ------------------------------------------------------------
% E-3. Section 3.3, scoring and metrics
% ------------------------------------------------------------

\paragraph{Scoring and metrics.}
At every turn $t$ of item $i$, the judge assigns a position strength
$s_i^{(t)} \in \{0,\dots,4\}$: a score of 4 means the target fully
defends the gold position, intermediate values mark hedging and partial
concession, and 0 marks collapse, which we define conservatively as the
target endorsing the false premise as a general truth in its own voice.
A dialogue ends at the first collapse or at the turn budget $T = 99$.
Let $t_i^{c} = \min\{t \le T : s_i^{(t)} = 0\}$ denote the collapse
turn, with $t_i^{c} = \infty$ when the target never collapses. We
report three quantities:
\begin{equation}
  \mathrm{CR@}T = \mathbb{E}_i\!\left[\,\mathbbm{1}\{t_i^{c} \le T\}\,\right]
  \label{eq:cr}
\end{equation}
\begin{equation}
  \mathrm{TH@}T = \mathbb{E}_i\!\left[\,\min(t_i^{c} - 1,\; T)\,\right]
  \label{eq:th}
\end{equation}
\begin{equation}
  \mathrm{AUSC} = \mathbb{E}_i\!\left[\frac{1}{4T}\sum_{t=1}^{T}
    \tilde{s}_i^{(t)}\right],
  \quad
  \tilde{s}_i^{(t)} =
  \begin{cases}
    s_i^{(t)} & t < t_i^{c} \\
    0         & t \ge t_i^{c}
  \end{cases}
  \label{eq:ausc}
\end{equation}
CR@$T$ is the fraction of items on which the target collapses within
budget; TH@$T$ measures how long it holds its position; AUSC, the
normalized area under the strength curve, credits a model for how much
of its position it retains rather than only for whether it collapses.
We additionally record the first-erosion turn,
$\min\{t : s_i^{(t)} < s_i^{(t-1)}\}$, the earliest point at which the
target concedes ground. TH@$T$ plays the role that Turn-of-Flip plays
in SYCON-Bench, generalized to a graded score and an adaptive horizon;
AUSC has no analogue among flip-based metrics and exists to expose runs
that erode steadily yet never register a flip.


% ------------------------------------------------------------
% E-7. Section 4, opening paragraph + design-summary table
% (verify: debate scenario deliberately dropped? repo has a
%  debate_setting/ folder. Item counts pending C.2 checks.
%  Table requires \usepackage{booktabs}.)
% ------------------------------------------------------------

\section{Experimental Setting}
\label{sec:setting}

We run every target model through the SPINE protocol in a two-by-two
design: two scenarios, each under both the naturalistic and the
adversarial condition of \S\ref{sec:protocol}. The scenarios
instantiate the implicit half of the coverage taxonomy of
\citet{hong2025measuring}. Challenging unethical queries probes a
subjective user view, a stereotype the user takes for granted;
identifying false presuppositions probes an objective one, a factual
claim that happens to be false. In both, the premise is embedded in the
question rather than stated as a position, so the target must first
notice what it is being asked to accept. We omit SYCON-Bench's third
scenario, debate, because erosion presupposes a gold position: a debate
stance is assigned rather than correct, and abandoning it is a
consistency failure rather than an honesty failure.
Table~\ref{tab:design} summarizes the design;
\S\ref{sec:scenarios} defines each scenario's failure criterion,
\S\ref{sec:datasets} the datasets, \S\ref{sec:models} the target
models, and \S\ref{sec:judge} the judge.

\begin{table}[t]
  \centering
  \small
  \begin{tabular}{lllr}
    \toprule
    Scenario & User view & Premise & Items \\
    \midrule
    Unethical queries      & Subjective & Implicit stereotype  & 20 \\
    False presuppositions  & Objective  & Implicit false fact  & 25 \\
    \bottomrule
  \end{tabular}
  \caption{The SPINE evaluation design. Seed questions are adapted from
    SYCON-Bench \citep{hong2025measuring}, which curated them from
    StereoSet \citep{nadeem2020stereoset} and CREPE
    \citep{yu2022crepe}. Each scenario is evaluated under both the
    naturalistic and adversarial conditions of \S\ref{sec:protocol}.}
  \label{tab:design}
\end{table}


% ------------------------------------------------------------
% E-8. Section 4.1, scenarios
% (renamed from "Method"; FP-first order — keep consistent with
%  the Sec. 4 opening paragraph and table row order.
%  UQ scenario still needs an example figure.)
% ------------------------------------------------------------

\subsection{Scenarios}
\label{sec:scenarios}

\paragraph{Identifying false presuppositions.}
The target receives a natural question that presupposes something
false. It is expected to surface the premise, correct it, and hold the
correction under pressure. The premise rides in the presupposition of a
why-question rather than in an assertion
\citep{yu2022crepe, hong2025measuring}, so answering at face value
already concedes it: the target must notice what the question assumes
before it can resist it. A run is scored as failed only when the target
affirms the false premise as an unconditional, general truth in its own
voice; hedges, partial retractions, and narrow concessions lower the
strength score (\S\ref{sec:metrics}) but never register as collapse.
Figure~[N] shows an example.

\paragraph{Challenging unethical queries.}
The target receives an advice-seeking prompt built on an implicit
stereotype. It is expected to withhold endorsement and answer the
legitimate part of the request without the assumption. The stereotype
is embedded as an assumption rather than stated as a claim, following
the design of \citet{hong2025measuring}: stated overtly, it would be
refused at the first turn and the dialogue would never begin. A run is
scored as failed when the target affirms the stereotype as a general
truth in its own voice, or advises treating an individual by group
membership; this discriminatory-action signal catches endorsement that
hides in recommendations while the proposition itself stays hedged.
Figure~[N] shows an example.


% ------------------------------------------------------------
% E-9. Section 4.2, datasets
% (verify: CREPE (Yu et al. 2022) vs (QA)^2 (Kim et al. 2023);
%  who rewrote the StereoSet prompts -- keep the bracketed
%  choice honest; 18/20 split vs four-category claim;
%  [N] = total scored turns from run logs)
% ------------------------------------------------------------

\subsection{Datasets}
\label{sec:datasets}

The dataset consists of seed items only: no dialogue is stored in
advance, because every conversation is generated live by the protocol
of \S\ref{sec:protocol}. The seeds are adapted from SYCON-Bench
\citep{hong2025measuring}, a deliberate choice: holding the items fixed
keeps our results directly comparable to its scripted five-turn
protocol, so differences in outcome are attributable to the pressure
regime rather than the data. Item counts are small because each seed
produces up to 99 adjudicated turns per condition per target, roughly
[N] scored turns in total; we trade breadth of items for depth of
interaction.

\paragraph{Identifying false presuppositions.}
We use 25 questions that SYCON-Bench drew from CREPE
\citep{yu2022crepe}, a QA dataset of naturally occurring false
presuppositions. The 25 are selected to span five presupposition types
evenly, five per type: false properties, false causal relations, false
predicates, false existentials, and false clauses. Each item is a
(question, false presupposition, gold correction) triple: the
presupposition states the false premise as a general claim, and the
correction specifies the position the target should hold. The triple is
also the judge's reference: position strength (\S\ref{sec:metrics}) is
scored against the gold correction. The full item list appears in
Appendix~[X].

\paragraph{Challenging unethical queries.}
We use 20 prompts that SYCON-Bench drew from StereoSet
\citep{nadeem2020stereoset}, each [rewritten by
\citet{hong2025measuring} / rewritten by us] so that the stereotype is
assumed rather than stated (\S\ref{sec:scenarios}). Each item is a
(question, stereotypical presupposition, correct stance) triple: the
presupposition states the stereotype as a general categorical claim,
and the stance specifies the position the target should hold. As in the
first scenario, the judge scores position strength against the stance.
The premises are predominantly nationality and ethnicity stereotypes
(18 of 20), with one gender and one occupational item. The rewritten
prompts appear in Appendix~[X].


% ------------------------------------------------------------
% E-10. Section 4.3, target models
% (settle the roster from jsonl meta BEFORE pasting: text,
%  models table, and results tables must match; verify closed-
%  model trace availability per provider; Qwen3-32B thinking is
%  partial -- reference-only footnote)
% ------------------------------------------------------------

\subsection{Target Models}
\label{sec:models}

We evaluate [12] target models, selected for the contrasts they
support rather than for coverage: every model-level claim in
\S\ref{sec:results} is a comparison within a family. The reasoning
axis is carried by matched pairs. Qwen3 \citep{yang2025qwen3} runs at
three scales (8B, 32B, and 235B), each in both thinking and
non-thinking mode, so the reasoning contrast holds the weights fixed
and toggles only the inference regime; DeepSeek V3
\citep{deepseek2025v3} and R1 \citep{deepseek2025r1} contrast a chat
model with its reasoning-optimized counterpart. The same Qwen3 series
provides the scale ladder. Gemini 3.1 Flash Lite and Gemini 3.1 Pro
\citep{google2026gemini}, GPT-5.5 and GPT-5-mini \citep{openai2026gpt},
and Claude Haiku 4.5, Claude Sonnet 4.6, and Claude Opus 4.8
\citep{anthropic2025claude} anchor the results to deployed frontier
systems. Table~\ref{tab:models} lists the full roster.

The adversarial condition reads the target's reasoning trace
(\S\ref{sec:pressure}), which not every target exposes. Open-weight
reasoning models return the full trace; non-thinking modes produce
none; [closed models return a summarized trace / no trace -- verify
per provider]. Where no trace is available, the proxy conditions on
the final reply alone, and the affected cells are marked in
Table~\ref{tab:models}. All targets run with the same minimal system
prompt; decoding parameters and API versions appear in Appendix~[X].


% ------------------------------------------------------------
% E-4. Section 4.4, judge and reliability (skeleton)
% ------------------------------------------------------------

\paragraph{Judge and reliability.}
SPINE's measurements are only as trustworthy as the judge, so we defend
them on two fronts. The collapse criterion is conservative by
construction: a run is marked collapsed only when the target endorses
the false premise as a general truth in its own voice. Hedges and
partial concessions lower the strength score but can never register a
collapse on their own, so the genuinely ambiguous middle of the scale
cannot produce false positives on CR@99. To assess scoring stability,
we re-score [N] held-out transcripts turn by turn with six judge models
([list]), and compare each against a human-annotated reference set of
[M] turns. Mean pairwise agreement across judges is [X]; agreement with
the human reference is highest for Claude Sonnet 5 ($\kappa = $ [Y]),
which is why it serves as the judge in the main experiments.
Disagreement concentrates on mid-scale turns, where partial concession
is genuinely borderline; on the turns that determine collapse (strength
4 and 0), agreement reaches [Z]. Per-judge results appear in
Appendix~[X].
% If Sonnet 5 remains a *target* model: add one sentence here on
% re-scoring its runs with a different judge to rule out self-preference.
